% Source: https://artofproblemsolving.com/community/c3h2445046_elementary_probability_problem
% Stable identity: aops-topic-2445046-q01
The number of zeros at the end of $m!$ is the number of factors $10=2\cdot5$ that can be formed from its prime factorisation, i.e. the smaller of the exponents of $2$ and $5$ in $m!$. There are always at least as many factors $2$ as factors $5$ (among $1,\dots,m$ the multiples of $2$ are at least as frequent as those of $5$, and similarly for $4$ and $25$, and so on), so the number of zeros is the exponent of $5$ in $m!$.

Since $(k+1)!=(k+1)\cdot k!$, the number of trailing zeros goes up by exactly the exponent of $5$ in $k+1$. So we need $k+1$ to be divisible by $25$ but not by $125$. With $1\le k\le100$ we have $2\le k+1\le101$, so $k+1\in\{25,50,75,100\}$ and $k\in\{24,49,74,99\}$. The probability is
\[\dfrac{4}{100}=\boxed{\textbf{(D)}\ \dfrac{1}{25}}.\]
(For example, $24!$ ends in $4$ zeros and $25!$ in $6$.)
